
Semantic Role Labelling~\citep[SRL, ][]{marquez08srl} is generally understood as 
the task of identifying and classifying the semantic arguments and modifiers of 
the predicates mentioned in a sentence. For example, in the case of the 
following sentence\footnote{``Haag plays Elianti'' is a segment of a sentence in 
training corpus.}:
\begin{quote}
\begin{center}
    \includegraphics[scale=.63]{haag-example}
\end{center}
\end{quote}
we are to find out that for the predicate token {}``plays'' with sense ``play a 
role'' (play.02) the phrase headed by the token {}``Haag'' is referring to the 
player (A0) of the play event, and the phrase headed by the token {}``Elianti''  
is referring to the role (A1) being played. SRL is considered as a key task for 
applications that require to answer {}``Who'', {}``What'', {}``Where'', etc.  
questions, such as Information Extraction, Question Answering and Summarization.  

In this paper we present a Markov Logic \citep[ML,][]{richardson06mln} to define 
a multi-language joint SRL model. Section \ref{sec:introduction} introduces ML.  
At the core of ML is a set of First Order Logic (FOL) formulae which models the 
SRL task. We present our specification of the model in section \ref{sec:model}.  
ML allows to model the task independently of the language, formulae from one 
language can be used in different languages. Other formulae can be language 
specific.  

